% 译者说明
\section*{译者说明}
\addcontentsline{toc}{section}{译者说明}
\markboth{译者说明}{译者说明}

\subsection*{关于底本}

本书是 M. Kontsevich 与 Y. Soibelman 的论文 \emph{Homological mirror symmetry and torus
fibrations} 的中文译注本。底本为 arXiv:math/0011041v2 [math.SG]（2001 年 6 月 4 日提交，
共 66 页，文章编号沿用作者自排的 letter 版式）。底本 PDF 的扉页上另有一行排版日期
「November 26, 2024」，那是文献服务器后来重新编译时留下的时间戳，并非论文的写作或修订年份；
文中所有历史陈述（例如「在最近的一篇预印本 [GW] 中」）都应按 2001 年的语境理解。

底本没有摘要，正文自 §1 直接开始。全书按原文章节编排：§1 引言，§2 共形场论方面的物理动机，
§3 Calabi–Yau 流形的大复结构极限，§4 A∞-代数与 A∞-范畴，§5 Fukaya 范畴及其退化，
§6 Morse–Smale 复形与 Morse 函数范畴，§7 凝聚层导出范畴的 A∞-结构，§8 同调镜像猜想，
§9 复几何情形下的构造（附录）。

\subsection*{翻译体例}

\begin{enumerate}[leftmargin=2.2em,itemsep=2pt]
\item \textbf{定理编号照原文。}原文每一类环境各自连续编号，不冠节号：定理 1–4、命题 1–15、
      引理 1–3、推论 1、猜想 1–5、定义 1–23、注 1–22。译文沿用这一编号，读者可据此与原文直接对勘。
      原文中的「§6.3 中的定理」指的是 §6.3 一节里的定理，而不是「定理 6.3」。
\item \textbf{术语。}经典术语沿用中文学界通行译法，人名一律不音译。每个术语首次出现时以
      「中文（English）」的形式给出原文，其后只用中文。全书译名见附录 A 的术语表。
\item \textbf{译注。}凡属术语背景、原文省略之处、以及原文的笔误或可疑之处，均以橙色方框「译注」标出。
      译注只作说明，不改动原文的数学内容；原文有误时译文照原文译出，另加注说明。
\item \textbf{公式。}所有公式均据底本页面图像逐符号核校，未采用文字提取层的转写。
\item \textbf{参考文献。}原文 p63–66 的参考文献表保留英文原样，不作翻译，以便读者按原标签
      （\texttt{[SYZ]}、\texttt{[Ko]} 等）检索。
\end{enumerate}

\subsection*{致读者}

这篇论文写于 2001 年，其中若干设想（特别是 §8 关于阿贝尔簇情形的同调镜像对称、
以及 §9 的 Dolbeault 扇形分解）在随后二十年里获得了长足发展，术语与记号也有变化。
译者只在必要处加注说明后来的进展，不把这些后来的结果混入正文，以免造成「作者当年即已如此陈述」
的误解。凡译注中引述的后来文献，均单独标明，不与原文的参考文献混同。

\medskip
\noindent 译文中的一切错误由译者负责。读者若发现误译或排印错误，请对照原文页码指出，
以便再版时订正。原文页码与译本页码的对应关系见附录 B。

\vspace{1.2em}
\noindent\hfill 译者\quad 二〇二六年
